\chaptertitle{9}{结合估值的定投，让你赚钱更容易}

\subsection{为什么指数基金长期收益不错，但投资者不赚钱}
在前面几个章节中，我们介绍了很多不错的指数基金。有的朋友看完了这些章节，会有一些疑问：我看这些指数，动不动就10年上涨三四倍、五六倍，甚至涨幅高达10倍的都有。但是现实生活中，身边的同事、朋友，很多人买股票基金不仅不赚钱，反而是亏钱的。这到底是怎么回事呢？

{\heiti “七亏二平一赚”是什么意思}

“七亏二平一赚”，意思是10个投资股票基金的人，有7个是亏损的，2个是不赔不赚的，1个是有盈利的。

但是国内的股票基金的净值，整体长期是上涨的。2018年5月16日，证监会专门写了一篇文章，里面提供了来自官方的基金收益数据。原文如下。

{\kaishu 截至2017年年底，公募基金偏股基金年化收益率平均为16.5\%，债券基金平均年化收益率为7.2\%。公募基金行业累计分红1.71万亿元。为投资者整体创造了很不错的长期价值。}

16.5\%，这不是总收益，而是每年平均的年化收益率！

这已经是相当可观的一个收益率了，超过我们普通投资者日常生活中常见的大多数投资品种。

股票基金这个品种在国内的收益率本身很不错，这一点毋庸置疑。

但是，很多股票基金的投资者并未赚到钱。

之前有一个网站针对投资者做过回测：仅仅是持有现金，可以跑赢接近70\%的股票投资者。换句话说，有70\%的股民投资股票的收益还不如只持有现金。如果简单买一个银行理财产品，能跑赢股票投资的比例还会更高。

这个比例放在股票基金上，也是成立的。很多人投资股票基金，也是被套的。

{\heiti 追涨杀跌导致投资者的收益不理想}

为什么股票基金整体是赚钱的，但是股票基金投资者是亏损的呢？

一个很重要的原因，就是这些投资者往往是在股市大幅上涨之后开始入市，并且往往是涨得越多，买得也越多。最后买在了高位，之后市场下跌，收益自然是惨淡的。

这也是新手投资的一个缺陷：追涨杀跌。

我们看看2005—2015年，国内所有新增基金开户数的走势，再对比一下同期的A股走势，如图9.1所示。

大家发现了什么呢？

（1）两种走势高度重合。在市场比较低迷的阶段，基金开户数也比较低迷。很少人在熊市对股票基金感兴趣，但是通常熊市是投资指数基金的好时机。

（2）在市场短期大涨指数基金变得比较贵之后，才会有大量的人开通账户，开始投资股票基金，例如2007年牛市和2015年牛市。

\begin{figure}[htbp]
\centering
\includegraphics[width=\textwidth,height=0.82\textheight,keepaspectratio]{figure/Image00031.png}
\end{figure}
图9.1 2005—2015年，国内所有新增基金开户数与同期A股走势

在过去的十几年里，有超过70\%的投资者，是在牛市的阶段才开始投资股票基金。市场涨得越多，股市的财富效应越明显，对投资者的吸引力越大，进入市场的投资者也越多。

当市场上涨到5000～6000点的时候，也是进入市场的股民数量最多的阶段。这些投资者在便宜的时候没有买入，在贵的时候买入了，自然就会亏钱。

投资者投资指数基金最大的一个误区，就是认为一只指数基金好就自始至终投资它，不管贵贱。

实际上，仔细想一想就知道这并不合理。

超市买西瓜，一元一斤还可以接受，但忽然有一天涨到30元一斤，我们就不会去买了。

指数基金也是如此，同样的指数基金，如果变贵了，相比较之前，它的投资价值就会下降。所以指数基金并不是一直都处于适合定投的状态的。

我们来看看王大爷的例子。

{\kaishu 王大爷在2014年7月的时候，买了2万元的沪深300指数基金，到了2015年的5月，王大爷打开账户一看，发现当初买的2万元的基金变成4万元了。不到一年的时间里，收益直接翻番，相当惊人。王大爷很高兴啊，这只基金这么好，这么容易赚钱，当初应该多买点。}

{\kaishu 于是王大爷在2015年5月又拿出10万元，买了同样一只基金。结果，这一次投入的10万元，到了2016年年初，不仅一分钱没赚，还赔了一万多。}

{\kaishu 王大爷就很纳闷，为什么同样一只指数基金，第一次买能赚钱，第二次买就亏了呢？}

{\kaishu 原因很简单，第二次，王大爷买贵了。}

买指数基金和买房子是一个道理。价值100万元的房子，你花200万元去买就买贵了。房子买贵了，再卖出去说不定会亏本；同样，指数基金买贵了，我们再想赚钱就很困难了。所以，我们要在指数基金便宜的时候去买，而不是贵的时候去买。

这也就是我们对定投指数基金的另一个优化：指数基金有自己的投资价值，我们要在指数基金值得投资的阶段买入。如果在牛市大涨后投资指数基金，反而容易遭遇亏损。

通常来说，指数基金在熊市值得坚持定投，而短期大涨之后进入牛市，指数基金包括其他股票基金，反而不适合投资。

那么问题来了，我们怎么判断是熊市还是牛市，以及适不适合开始投资指数基金呢？

这个问题的答案，就是估值。

\subsection{什么是估值}
{\heiti 没有哪种估值方法，可以精准计算价值}

简单来说，估值就是估算一个投资品种的价值，知道这个品种大概值多少钱。

类似的概念，其实我们在日常生活中都在使用。例如，我们去超市买西瓜，如果面前摆着一个西瓜，并没有标价格，我们称了一下是10斤重，那这个西瓜多少钱呢？我们凭经验估算，夏天通常一个西瓜的单价是0.6～0.8元/斤，那10斤重的西瓜，应该是6～8元。这个过程，就是给西瓜的一次估值。

我们能不能准确地估算出这个西瓜的价格呢？并不能，可能这个超市的西瓜是卖0.67元/斤。我们靠估算，是没有办法得出准确的数字的。但是我们可以给这个西瓜估值，确定在0.6～0.8元/斤这个范围。

因此，我们可以看出估值的第一个特点：没有哪种估值方法，可以精准地计算出价值。

我们估算出这个西瓜价格应该是0.6～0.8元/斤。结果问了一下超市的店员，这个西瓜目前价格是0.5元/斤，比我们估算的价格低不少，这是一个好价格。用这个价格买下这个瓜，我们不亏。

从这里，我们可以看出应用估值的方法：

在价格低于估算出来的价值时，买入，这样才划算。没有人愿意用1元买价值0.8元的东西。

这也是巴菲特的老师——本杰明·格雷厄姆（Benjamin Graham）提倡的投资方法：价值投资。估值是投资中很重要的一部分，学会估值方法，是实践价值投资的第一步。

{\heiti 估值是一门艺术，不是科学}

不过，不同的人，得出的估值结论并不一致。

继续使用上述例子，我们根据西瓜过去一年夏天的平均价格，估算出今年西瓜平均价格的合理范围，可能是0.6～0.8元/斤。但是，如果我们问一下种西瓜的瓜农，他们说今年种西瓜，肥料、水电、物流、人工的成本都比去年要高一些，要想不亏本，西瓜得卖0.9元/斤。瓜农认为，这个才是西瓜的合理价格，因为低于这个价格大多数瓜农会亏本。

瓜农说的有没有道理呢？也有道理。

估值并不是一个有明确客观标准的计算流程。估值本身是有一定主观性的。甚至在很多时候，都是公说公有理、婆说婆有理。更何况每个人都有自己的信息来源、知识储备和认知水平，同一件事情，不同人得出的结论是不一样的。甚至同一件事情，同一个人从不同的角度，得出的结论也不一样。

估值并不是一门严谨量化的科学。有句话说得好，估值是一门艺术而不是科学。

所以，当我们看到不同人用同一种估值方法得出不同结论的时候，不要惊讶。这才是正常的情况。

那是不是说我们就没有办法利用估值了呢？也不是。

格雷厄姆就提出了一个理论：安全边际。简单理解，我们可以给计算出来的估值，保留足够的安全边际。这样，即使估算有偏差，保留了安全边际之后也可以使用了。

例如，我们估算出西瓜的价值在0.6～0.8元/斤，那可能不是很保险。我们给估算出来的价格打个折扣，例如0.4～0.6元/斤。打折后就更低了，保留了更大的安全边际。

这样就把不稳定的、有偏差的数据，变成了一个可以依赖和使用的数据了。这就是安全边际的重要性。

\subsection{相对估值法和绝对估值法}
估值的方法有哪些呢？目前市面上流行的估值方法有几十种。但是，这么多的估值方法从原理上来说，可以分为两类：相对估值法和绝对估值法。

{\heiti 什么是相对估值法}

比如说卖房子，通常用的就是相对估值法。

例如，王大爷家里有两套房子，想卖掉一套。

卖房子自己得有一个的心理价位，王大爷就去房地产中介那里询问小区的房价。中介告诉王大爷，小区同户型的房子，这周卖了200万元。于是，王大爷就把这个作为自己房子的估值，心想有人来买他房子的时候，他开价不会比这个价格低。

第二个月，中介告诉王大爷，同小区的房子从200万元涨到210万元了，王大爷也同步提升了自己房子的出售价格。

王大爷通过比较的方法，完成了给自己房子的估值，这就是相对估值法的一种应用。

相对估值法，核心在于一定是有参照物的。这个参照物，可以是同类比较，可以是历史比较。

{\heiti 什么是绝对估值法}

我们可以看一个案例。

王大妈从楼下的银行买了1万元的3年期封闭理财产品，到期后总收益率估计为10\%。换句话说，3年后王大妈能一次性拿回1万元本金，外加1000元的收益。

但是，王大妈买完的第二天，王大爷生病了急需用钱。王大妈没办法，只能把邻居叫过来，对大家说，“我手里这个理财产品刚买的，谁要就出个价吧”。

邻居小明心想，我最近在投资指数基金，估计每年10\%以上的收益率问题不大。如果王大妈出的价格低，我买了，能赚到接近指数基金的收益，那也是不错的选择。这个理财产品3年后能拿回11000元的“本金+收益”，如果我想要赚每年10\%的收益率，那我现在的出价应该是11000/1.1\textsuperscript{3} =8264元。所以，只要我出价不高于8264元，就可以从王大妈的这个理财产品上，赚到每年超过10\%的年化收益率了。

小明的这个估值过程，就是绝对估值法。

绝对估值法的核心，是围绕投资品种未来所能产生的现金流来的。

一个投资品种在它生命周期内所能产生的现金流的折现值之和，就是这个品种的绝对估值。

{\heiti 两种估值方法的应用}

在现实生活中，这两种估值方法都经常能遇到。

实际上无论用哪一种方法给一个品种估值，最后的结论应该是差不多的。如果一个品种处于低估便宜的阶段，那用相对估值法和绝对估值法，都可以得出“这个品种低估”的结论。

通常一个资产的未来现金流越明确，绝对估值法计算得就越精准。但是股票等资产的未来现金流比较难精准预测，所以绝对估值法在股票品种的估值上，经常会有很大偏差。从原理上来说，绝对估值法是更贴近一个品种的真实价值的。

实际应用的时候，这两种方法我们都可能会用到。结合起来，可以帮助我们更好地完成估值。

\subsection{在估值过程中，我们要关注哪些因素}
{\heiti “奥卡姆剃刀”原理}

影响估值的因素有很多，比如说今天市场又出了什么利好消息、明天又有哪家公司的业绩大涨。影响股票涨跌的因素，可能有数千种，实际上我们不可能同时兼顾太多的因素。

14世纪的英国逻辑学家奥卡姆（Ockham）提出一个原理：“如无必要，勿增实体”。这个原理被称为“奥卡姆剃刀”原理。

他认为，事情应该从复杂变为简单。如果解决一个问题的方法有很多种，我们应该选择最简单的那一种。如果一个问题有多种解释，那需要最少步骤和假设的解释更有可能是正确的。

对于“如何估值”来说，就是：当一个问题有多种解决方案的时候，那就应该选择最简单的解决方案。这个解决方案包括最基本的重要元素即可。如果能用3个指标就可以完成估值，那用6个就是自找麻烦。

{\heiti 约翰·博格：对指数基金收益影响最大的三个因素}

“指数基金之父”约翰·博格（John Bogle）在分析指数基金的时候，就应用过“奥卡姆剃刀”原理。

影响指数收益的因素实在是太多了：上市公司的盈利、国家的政策法规、经济环境、企业的竞争、管理层的好坏、投资者的心理情况、资金面、汇率波动……每一个因素都有可能在短期里影响指数的走势。

不过约翰·博格应用“奥卡姆剃刀”原理，总结出对指数基金收益影响最大的三个因素为市盈率、盈利、分红。通常明白这三个因素的变化，我们就能判断出指数基金的长期收益。另外，企业的净资产也是很重要的一个因素。

这就是我们给指数基金估值，最常用的几个因素。我们整套估值体系，也是建立在对这些因素的分析上，这样就避免了过多因素对我们估值分析的干扰。

对指数基金投资者来说，指数基金是一个简约版的股票资产。一家上市公司该有的净资产、盈利等，指数基金作为一篮子上市公司股票的集合，也是都有的。而且，相对于个股来说，指数基金的盈利、净资产、分红会更加稳定，因此在给指数基金估值的时候也更容易、成功率更高。

所以同样的估值方法应用在指数基金上，会更适合普通投资者。

这是对指数基金进行估值的优势。下面我们就来了解一下，这几种估值指标是什么。

\subsection{估值指标一把手——市盈率}
市盈率指标出现很长时间了，它使用方便、含义清晰，可以说是应用范围最广泛的估值指标了。

不过，市盈率指标也有一些缺点，如不适合用在盈利变化大的品种上，不适合强周期性品种，遇到金融危机和经济周期的底部会失效等。

所以也有很多金融学家在不断优化市盈率。

{\heiti 静态市盈率}

静态市盈率是指，市值/上一年度财报的净利润。比如说在2019年年中，我们看静态市盈率，就是用市值/2018年财报的净利润。

{\heiti 滚动市盈率}

时效性更高的滚动市盈率是指，市值/最近四个季度财报的净利润。相对来说实效性高一些。

螺丝钉经常用的就是滚动市盈率。

滚动市盈率每个季度更新一次。比如说在2019年7、8月，上市公司会更新半年报。半年报更新完成后，上市公司如果盈利上涨，滚动市盈率就会下降。

静态市盈率和滚动市盈率都是根据已经公布的财报来计算的，是客观的估值指标。只不过时效性不同，相对来说，滚动市盈率的时效性会更强一些。

{\heiti 动态市盈率}

预测未来的动态市盈率是指，市值/预测未来一年的净利润。动态市盈率通常是用未来一年的盈利来计算估值，但是这个未来一年的盈利，是预测的。

未来谁也无法准确预测，所以就会有偏差。大多数时候，我们都不用动态市盈率。

{\heiti 席勒市盈率}

罗伯特·席勒（Robert Shiller）是著名的经济学家，《非理性繁荣》（\textit{Irrational Exuberance}）就是他的著作。这里的席勒市盈率，是他对市盈率的一个改良。

一般的市盈率，使用一年的净利润时会遇到一些特殊情况而产生较大波动。所以巴菲特的老师格雷厄姆建议用多年的平均净利润来作为参考。

席勒在这个基础上，编制了席勒市盈率，也被称为CAPE。席勒市盈率使用过去10年的平均净利润作为分母，也就是市值/过去10年平均净利润。

这种方式的市盈率的好处，是避免了某一两年特殊情况的干扰。比如说2007—2008年金融危机的时候，美国上市公司整体盈利大幅下滑，标普500指数的静态市盈率被动上升到了上百倍。但是如果用席勒市盈率，则避免了这种干扰，继续有效。

通常来说，盈利是长期上涨的，所以10年的平均净利润，会比最近一年的净利润要低。所以，席勒市盈率比静态市盈率要高一些。

需要注意的是，席勒市盈率是用来衡量一个股票市场整体情况的，特别是对有较长历史的市场比较适合。比如说国内的创业板，历史太短，就用不了席勒市盈率，但像历史比较悠久的港股市场是可以的。

{\heiti 中位数市盈率}

中位数市盈率，是指数中才有的一个估值指标，个股中反而没有。

所谓的中位数市盈率，就是将一个指数的成份股，按照市盈率从高到低排序，取排在中间的。

比如说上证50指数有50只成份股，市盈率中位数就是排中间位置的第25位股票的市盈率数值。

中位数市盈率有什么用呢？

有的指数，部分成份股占据指数绝大多数比例，这样的指数，市盈率可能会被部分权重股票“绑架”。

什么意思呢？

就是说指数中，有可能其他股票已经大涨，但权重股票还没有大涨，显得指数“还处于低位”。这时，中位数市盈率就派上用场了。

中位数市盈率会更敏感，在牛市的时候，会比指数本身的市盈率波动更大。这样避免了权重股“绑架”指数导致估值失真的影响。

那什么样的指数适合参考中位数市盈率呢？

如果指数成份股的平均市值规模差别并不大，那这样的指数，参考中位数市盈率意义不大。

如果指数成份股的平均市值规模差别比较大，那这样的指数，可以辅助参考一下中位数市盈率。

例如牛市的时候，如果中位数市盈率先进入了高估区域，那说明这个指数其实大部分股票都涨了比较多，可能有风险了。

各种不同的估值指标，都是我们的武器。谈不上哪个指标更先进。从多个角度考虑，能更好地为我们的投资保驾护航。

\subsection{估值指标二把手——市净率}
除了市盈率，市净率也是我们常用的估值指标。市净率通常是作为二号指标来参考的。

{\heiti 什么是市净率}

市净率是指每股股价与每股净资产的比率，也就是我们说的账面价值。

市净率的定义是：PB=P/B（其中P代表公司市值，B代表公司净资产）。

净资产，通俗来说就是资产减去负债，它代表全体股东共同享有的权益。具体的计算在上市公司的年报中都有。

净资产这个财务指标比盈利更加稳定。而且，大多数公司的净资产都是稳定增加的，可以计算出市净率。

相比较市盈率来说，市净率会更稳定。因为盈利是一个时间段里产生的，而净资产是积累起来的。

某一年的企业经营好坏可能会对企业的盈利产生比较大的影响，但比较难对企业的净资产产生太大的冲击。

或者说，只要一家企业还没有亏损，大多数时候，净资产都是增长的，市净率就可以参考。

{\fangsong 影响市净率的因素}

{\heiti 企业运作资产的效率：ROE}

说到市净率，就不得不说净资产收益率（ROE），ROE等于净利润除以净资产。

对企业来说，资产本身是企业运营的素材，企业需要运作资产，并由此产生收益，资产才有价值。

针对同样的资产，有的企业可以运作它，产生更高的收益，获得更高的资产运作效率。衡量资产运作效率的重要指标就是ROE。

个人认为，ROE是企业所有经营指标中最关键的一个。

巴菲特的老搭档查理·芒格（Charlie Munger）也说过：

{\kaishu 从长期来看，一只股票的回报率与企业发展息息相关，如果一家企业40年来的盈利一直是资本的6\%，即ROE为6\%，那40年后你的年均收益率不会和6\%有什么区别，即使你当初买的是便宜货。}

{\kaishu 如果该企业在20～30年的盈利都是资本的18\%，即使当初出价高，其回报依然会让你满意。}

ROE越高的企业，资产运作效率越高，市净率也就越高。

{\heiti 资产的价值稳定性}

资产的种类多种多样，有的资产能够随着时间增值，有的则会随着时间迅速贬值。

例如茅台，正在酿造中的白酒会随着时间不断增值，但是英特尔公司生产的电脑芯片，如果没有及时销售，就会迅速贬值。

资产价值越稳定，市净率的有效性越高。

{\heiti 无形资产}

传统企业的净资产很大一部分是有形资产，例如地皮、矿山、厂房、原材料等。它们的价值比较容易衡量。

但是也有很多资产是无形资产，其价值很难衡量。例如企业品牌、高级技术工程师、企业专利、企业的渠道影响力、行业话语权等。

如果一家企业主要靠无形资产来经营，像律师事务所、广告服务公司、互联网公司，那市净率就没有多少参考价值了。

{\heiti 负债大增或亏损}

净资产是公司的资产减去负债。

如果公司的负债不稳定，极有可能会干扰净资产。公司如果出现亏损，可能会侵蚀更多的资产，导致净资产减少。

从这里我们可以看出，当企业的资产大多是比较容易衡量价值的有形资产，并且是长期保值的资产时，比较适合用市净率来估值。

还有一点是，行业要始终都是赢利的。比如说证券行业，在熊市的时候盈利是大幅下滑的，不过即便是大幅下滑，大多数时候证券行业也是挣钱的。这样净资产就不会受到太大的侵蚀。

但是如果在周期底部，企业经营不好，整个行业大范围亏损，那这样的行业，就不适合参考市净率了。

另外，在一些比较特殊的时间段，如遇到短期的经济危机，原本盈利稳定的宽基指数基金会在短期内盈利不稳定，也可以用市净率来辅助估值。

像2008—2009年美国次贷危机的时候，美股的盈利就出现了短期的大幅下滑。

这种下滑是不会长期持续的，但确实会影响用市盈率估值。这时，就可以用市净率来辅助进行估值。

之前有朋友问，如果盈利下滑导致市盈率被动上涨，这种情况下市盈率不能参考了，该怎么办。

此时，我们就可以用二号估值指标——市净率。

\subsection{衡量现金分红收益率的指标——股息率}
{\heiti 什么是股息率}

股息率是指，现金分红/市值。

股息，是投资者在不减少所持有的股权资产的前提下，仍然可以直接分享企业业绩增长的最佳方式。

我们都知道，有的股票是会分红的。如果我们用过去一年公司的现金分红除以公司的市值，得到的就是股息率了。

股息率衡量的就是现金分红的收益率。

那么，每年的现金分红从哪里来呢？这主要就是来自企业的盈利了。

企业每年会拿出一部分盈利，以现金分红的形式，分给股东。

股票，很早以前就出现了。但是，股票市场是最近两三百年才兴起的。在股票市场出现之前，人们交易股票，主要是自己私底下交易，很不方便。所以，当时绝大多数股票持有者都是长期持有股票，依靠股票的分红来获取收益。

{\heiti 分红比例是什么}

说到股息率，就不得不说分红比例，也就是分红率。乍一看，股息率、分红率这两个概念很像，但实际上是不同的。分红率，也叫作派息率、分红比例。

简单来说，股息率，是企业过去一年的现金派息额除以公司的总市值。分红率，则是企业过去一年的现金派息额除以公司的总净利润。

这两个比率，分子一样，但分母是不同的。一个是公司总市值，一个是公司当年总净利润（当然，细节上还有一些不同，例如分红还要扣税等，但大体上可以这么理解）。

分红率一般是公司预先就设置好了，并且连续多年都不会改变。例如中国工商银行等，分红率都在50\%。也就是说，中国工商银行会把过去一年净利润的50\%，以现金的形式进行派息。

而股息率，则会随着股价波动：股价越低，股息率就越高。

{\heiti 股息有什么用}

作为普通的二级市场股东，我们买入股票是享受不到很多权利的。

巴菲特可以用巨额资金将整家企业买下来。例如巴菲特用2500万美元买下喜诗糖果，而当年喜诗糖果的净利润有208万美元。

作为公司的实际掌控者，巴菲特可以自由抉择这208万美元的净利润是用来再投资还是花掉。但是，普通投资者是无法插手公司净利润的。

也就是说，即使公司的净利润高速增长，但是如果二级市场的股价不上涨，散户还是无法享受公司的业绩上涨带来的好处。

然而，股息可以将公司的净利润增长映射到投资者的现金收益上。

这是因为，股息是净利润的一部分，净利润上涨也会带动股息长期上涨。

你无须出售手里的股权，就可以获取不断增长的现金流。

实际上，这也是国家分享国有上市企业收益的方式：像国内的上市银行，大股东多为证金公司、财政部、发改委，它们要保证上市银行的国有控股，手里的股票很少会出售。

在这种情况下，享受收益的方法就是提高公司的分红率，从而每年获得高额股息。

股息是一个既能长期持有股权资产，又能同时享受现金流收益的好方法。

通过长时间持有高股息的股票资产组合，你可以无须关注股价涨跌，通过收获越来越高的现金分红来轻松实现财务自由。

这就是长期投资股票类资产最好的收益方式。我们无须再为波动的股价忐忑不安，每年收取稳定增长的现金分红即可。

{\heiti 时间越长，分红对收益的贡献越高}

时间越长，分红在我们投资股票的收益中的占比就越大。

美国著名金融学家西格尔教授，在他的著作《股市长线法宝》中分享了相关研究结果。

西格尔对标普500指数从1871年到2012年的数据进行回测，发现股息分红是整个时期内股东收益的最重要来源。

{\kaishu 从1871年开始，股票实际收益率为6.48\%，其中包括股息分红收益率为4.4\%，资本利得收益率为1.99\%。}

所以，投资指数基金的时候，股息率也是一个很重要的指标。

\subsection{股息率的神奇之处}
股息率有一个非常神奇的地方，就是现金分红通常不受股价涨跌的影响。

{\heiti 分红的神奇之处}

现金分红通常是上市公司盈利的一部分，所以除非经济出现非常大的变动，上市公司整体的盈利出现了问题，否则大多数时候，现金分红都是慢慢上涨的。

指数的股息率，说的就是上市公司的现金分红/市值。也就是说，分子是随着盈利增长而长期增长的。

我们以长期分红的红利ETF为例，每份基金分红如下。

{\kaishu 2012年：0.035元。}

{\kaishu 2013年：0.041元。}

{\kaishu 2014年：0.059元。}

{\kaishu 2015年：0.080元。}

{\kaishu 2016年：0.050元。}

{\kaishu 2017年：0.091元。}

{\kaishu 2018年：0.109元。}

{\kaishu 2019年：0.098元。}

如图9.2所示，长期来看分红是增长的。2019年的分红已经是2012年分红的两倍多了。只要继续持有，后面分红还会继续增长。这就是现金分红的神奇之处。

\begin{figure}[htbp]
\centering
\includegraphics[width=\textwidth,height=0.82\textheight,keepaspectratio]{figure/Image00032.png}
\end{figure}
图9.2 2012—2019年红利ETF每份基金分红金额

不过，长期增长不等于每年连续增长。偶尔一两年盈利下降，也会导致分红出现下降的。

{\heiti 投资领域}

正因为上述这个特点，高股息率的投资品种，也常常被用在一些投资领域。

{\heiti 领域1：保险资金}

保险资金很喜欢高股息率的品种，高股息率的股票是保险机构最爱的资产之一。

这是因为保险机构每年都需要支出一些费用，如果在高分红品种处于低估的时候买入，长期持有，那每年的分红就可以支付保费支出了。

比如说中国平安首席投资官陈德贤就多次介绍过平安投资股票的逻辑：寻觅分红比较高、估值相对便宜的股票标的。

{\heiti 领域2：养老投资}

我们之前介绍过的养老定投，也是建立在高分红的指数基金的基础上。

假如到老年的时候，我们的家庭每年需要支出10万元。如果我们持有一篮子高分红指数基金，这些指数基金的分红每年超过10万元，那我们就可以放心了。

因为分红不太受股价涨跌的影响，不管股票市场涨涨跌跌，每年的分红都是比较稳健的。这些分红，可以用来应对生活的支出。

引言介绍过的人物，也是持有大量的指数基金，依靠指数基金每年的分红来支付开销，从而实现40岁前退休、环游世界的目标。

这种投资理念，在欧美等成熟市场非常常见。因为欧美市场的分红比例比较高，愿意长期分红的股票数量也比较多。

不过比较遗憾的是，国内距离这一步，目前还有一段路要走。

{\heiti 股息率的高低如何判断}

那么，股息率多高算高，多低算低呢？其实不同的品种，情况也不一样。

通常来说，如果长期平均股息率比银行储蓄利率要高，那就已经不错了。

国内银行的一年期储蓄利率是1.5\%左右。不过这是一个基准数字。大部分银行的一年期储蓄利率都比这个高，像部分城商行，一年期储蓄的利率是2.5\%～3\%。

一只指数基金的股息率比银行的一年期储蓄利率要高，就已经算不错了的。持有这只指数基金，每年的分红都比储蓄的收益要高。

对于红利指数基金来说，这个条件不难满足。红利指数基金大部分时候的股息率都在3\%以上。当市场整体比较低迷的时候，红利指数的股息率能到4\%。

历史上，A股红利指数的股息率超过5\%的时候，都是历史大底，在这段时间里投资，收益都非常不错。

市盈率、市净率、股息率，是最常用的三个估值指标。

除此之外，还有一些估值指标也经常被用到，但使用频率会低一些，比如说市销率、市现率等。这些指标，我们在给指数基金估值的时候用得不多，更多的是个股投资的时候使用。所以，在这本书里就不过多介绍了。

至此，我们已经了解了各种不同的估值指标。但是这些估值指标具体该如何应用呢？

\subsection{不同的估值方法}
{\heiti 5种类型的指数基金}

如果从盈利的角度去观察指数基金，我们可以找到5种类型的品种。

\begin{itemize}
\item 盈利保持不变。
\item 盈利长期稳定增长。
\item 盈利快速增长。
\item 盈利大起大落。
\item 盈利衰退。
\end{itemize}

盈利大起大落、盈利衰退的品种，投资难度太高，并不适合作为普通投资者的投资主力。

所以我们主要针对的就是前3类品种。

最容易分析的是第1类，也就是盈利保持不变的品种。

\subsection{盈利保持不变：看现金流收益率高低}
如果一个品种的盈利始终保持不变，实际上这就是一个债券类资产。比如说我们去银行买一只3年期年化收益率为3\%的国债。买入10000元，第一年可以收到300元利息，第二年收到300元利息，第三年收到300元利息及10000元本金。

假如说有一家公司，这个公司没有其他的业务，它的全部业务就是把工资的资产买成这只国债。那这家公司的盈利，也会保持每年不变。

对这类品种的估值，通常看现金流收益率的高低。如果现金流收益率让我们满意，就可以考虑。

但前提是，这类品种的安全性要好，不能中途倒闭或者爆雷，让我们的收益受到损失。

通常来说，这类品种可以参考债券基金的收益，它的收益跟债券基金品种的平均收益接近，不会有太大的差别。如果这类品种的收益远高于债券基金的平均收益，那就要小心，背后可能有风险，例如违约风险，或者未来盈利大幅下滑的风险。

其次，这类品种的现金流收益率，一定要比国债收益率高。因为国债的安全性是最高的，如果是股票类品种，安全性是不如国债的，这类品种必须提供更高的现金流收益率，才值得考虑。

这一点很容易理解。但在投资过程中，到了牛市，经常有一些品种，现金流收益率不到2\%。

因为，现金流收益率=现金流/市值。牛市很多品种市值大涨，导致这些品种的实际收益非常低，但投资者看中短期的赚钱效应，反而去哄抢这些品种。这就是不明智的行为。

盈利保持不变的品种，估值是最简单的，跟买理财产品、买债券差不多。但绝大部分指数基金，并不是这一类型，而是盈利稳定增长、盈利快速增长、盈利大起大落这三类。

\subsection{盈利稳定增长：盈利收益率法}
盈利收益率法是很多朋友开始指数基金投资所接触的第一个估值方法，它适合盈利稳定增长的指数基金品种。

{\heiti 什么是盈利收益率法}

盈利收益率是市盈率的变种。市盈率是用公司市值除以盈利，而盈利收益率则是用公司盈利除以公司市值。换句话说，盈利收益率约等于市盈率的倒数。

盈利收益率法最早是由巴菲特的老师格雷厄姆提出的。

盈利收益率法的使用方法如下。

\begin{itemize}
\item 当指数基金的盈利收益率大于10\%且大于国债收益率两倍以上的时候投资。
\item 当盈利收益率小于债券基金预期收益率（目前为6.4\%）的时候卖出。
\end{itemize}

盈利收益率法，其实是绝对估值法的一种应用。

{\heiti 盈利收益率法是一种绝对估值法}

绝对估值法，是把一个品种未来的现金流折现到现在，计算出所有现金流的现值，如果股价低于这个现值，说明当前的价格被低估了，如果低估很多，我们就可以考虑投资。

这就涉及一个概念——折现。什么是折现呢？

比如说现在有个朋友来问我们借10万元，打算借3年。

如果朋友说，现在借我10万元，3年后我还你10万元。我们心里会不舒服，心想这10万元我就算存银行，假如一年5\%的收益率，3年下来也有1.5万元的利息，为什么要白白借给你。

如果朋友说，现在借我10万元，3年后我还你11.5万元。我们还是会有些不舒服，因为同样的钱存在银行，3年下来也是1.5万元的利息收益。银行的信用比你个人的信用强多了，我借钱给你，还得担心你能不能还上。

如果朋友说，现在借我10万元，3年后我还你14万元。我们算了一下，3年后14万元，年化收益率超过10\%了。这个收益率还不错，可以借。

这就是货币的时间价值的体现。我们借钱的时候，是知道货币有时间价值的。

但是3年后的100元，相当于今天多少钱呢？这个具体换算关系如何？

这就涉及一个概念——折现率。

折现率就是按照一个百分比，把未来的现金流折算成当前值多少钱。

比如说折现率为10\%，那1年后的100元，就相当于现在的100/（1+10\%）=90.9元。2年后的100元，就相当于现在的100/（1+10\%）\textsuperscript{2} =82.6元。

换一个角度理解，折现率就是我们想要的年复合收益率。

比如说我们想要10\%的年复合收益率，那就可以用10\%折现到现在，算一个价格。用这个价格买入持有，就可以获得10\%的年复合收益率。

从这里我们可以看出，如果我们要求的年复合收益率越高，使用的折现率越高，那折算到现在的价格就越低。那对应资产出现这么低的价格的机会也会越少。

不同人给出的折现率通常是不一样的。那折现率有没有一个通用的参考标准呢？还真是有的。

折现率不仅仅是用来做投资，做项目、做工程的时候，也经常遇到折现率的概念。比如说针对社会平均折现率，国家给出过一个参考数字：8\%。这个8\%是整个社会的平均折现率。

很多项目、工程在开始建设前，都要评估内部收益率，也就是IRR收益率。如果IRR收益率低于8\%，项目可能不予通过。

当然，这个折现率是社会的平均折现率。如果项目或者所要投资的行业，风险更大，那折现率也可以更高。也就是说，国家认为，我们在做工程投资的时候，年复合收益率至少要8\%以上才值得投、值得做。

对股票投资来说，波动风险更大，所以使用的折现率通常为10\%～12\%。换句话说，我们投资股票资产时的预期年复合收益率，至少是这个数字。

货币有时间价值。我们要用折现率，把未来的钱折算成现在的价格。

{\heiti 盈利收益率法：折现率10\%的绝对估值法}

盈利收益率法第一条规则是盈利收益率大于10\%，这也是使用盈利收益率法买入指数基金的主要规则。

这其实是一个绝对估值法的简单应用。

我们先看一下永续年金的计算公式：

{\kaishu 永续年金现金流的现值=每年的现金流/折现率}

再来看看盈利收益率的计算公式：

{\kaishu 盈利收益率=盈利/市值}

换句话说：

{\kaishu 市值=盈利/盈利收益率}

\begin{figure}[htbp]
\centering
\includegraphics[width=\textwidth,height=0.82\textheight,keepaspectratio]{figure/Image00033.png}
\end{figure}
图9.3永续年金现金流公式和盈利收益率公式的对应关系

其中，盈利就是指数基金的“现金流”；市值也就是指数基金目前的“现值”。所以这样对照一下，盈利收益率就是折现计算公式里面的“折现率”了。

格雷厄姆要求盈利收益率大于10\%的时候投资，就是要求投资者的折现率要大于10\%。

这就是盈利收益率法的核心：折现率10\%的绝对估值法。

{\heiti 盈利的增长}

这种计算方法，适合的是未来现金流固定不变的品种。但实际上，我们投资的指数基金的盈利会稳定增长。盈利稳定增长的品种，应该使用永续增长年金模型。

假如指数基金背后的公司，一年赚100万元。用永续年金模型，我们算出来现值是：100万元/10\%=1000万元。用永续增长年金模型，假设长期增长率为5\%，那算出来现值就是：100万元/（折现率-增长率）=100万元/（10\%-5\%）=2000万元。

只要指数基金背后公司的盈利是长期增长的，而不是不变的，那么用永续增长年金模型计算出来的现值，就比永续年金模型更高。哪怕盈利只是每年增长1\%，也是应该更高的。

那为何格雷厄姆不用永续增长年金模型呢？

{\heiti 安全边际}

这里体现的就是格雷厄姆的安全边际理论。

在盈利收益率大于10\%的时候买入，即使是对于一个“盈利以后再也不上涨”的永续品种来说，也是划算的。

如果买入后，盈利永远不增长，保持不变，那就变成“盈利保持不变”的“债券”了。利息收益率10\%的债券，其实也是不错的一笔投资。

在这个基础上，指数基金背后公司的盈利是上涨的，那我们就能赚得更多。

格雷厄姆做投资的年代，是20世纪发生第一次、第二次世界大战的时候。当时很多上市公司都是朝不保夕，根本谈不上发展。上市公司的整体盈利增长几乎停滞，或者增长缓慢。一直到战争结束后才恢复增长。所以，格雷厄姆在估值的时候，所用的估值方法是非常谨慎的。

世界上大部分国家的股票市场，长期盈利都是上涨的。所以盈利收益率法，过去几十年，在大多数国家股票市场都是很有效的。

这就是盈利收益率法的原理。它本质上是一个折现率为10\%的绝对估值法。

{\heiti 定投策略}

了解了盈利收益率法，那具体到定投的时候该如何应用呢？

在当前的国内利率和基金的收益水平下，我们使用盈利收益率来定投指数基金的策略如下。

\begin{itemize}
\item 当盈利收益率大于10\%时，分批投资。
\item 盈利收益率小于10\%，但大于6.4\%时，坚定持有已经买入的基金份额。
\item 当盈利收益率小于6.4\%时，卖出基金。
\end{itemize}

卖出后可以投资其他低估指数基金或者债券基金资产，等待下次机会出现。

在定投的过程中，有的朋友有疑问：“如果我这个月发现盈利收益率大于10\%，定投了，下个月盈利收益率小于6.4\%了，要卖出，这样操作会不会很频繁呢？”

其实，一般情况下，盈利收益率在短时间内变化不大，不会出现这么极端的情况。所以，我们只需要配合定投，在每月定投的时候，查看一下盈利收益率，再根据盈利收益率当时所处的范围，选择相对应的操作，也就是继续定投、坚定持有或是分批卖出，就可以了。

{\heiti 盈利收益率法的局限性}

盈利收益率法也有它的局限性。盈利收益率的使用条件比较苛刻，只适合于流通性比较好、盈利比较稳定的品种。如果是盈利增长速度较快或者盈利波动比较大的指数基金，则不适合使用盈利收益率法。

目前国内适合盈利收益率法的品种，主要是上证红利、中证红利、上证50、基本面50、上证50AH优选、央视50、恒生、恒生中国企业等指数的指数基金。这几个品种的投资很简单，当盈利收益率大于10\%时就可以投资，小于6.4\%时就可以卖出。

我们了解了盈利收益率法非常方便好用。不过，在投资指数基金的时候，并不只有盈利收益率法一种方法，而且由于盈利收益率法自身的局限性，只有有限的一些品种适合使用盈利收益率法进行投资。除了这些适合盈利收益率法的基金品种，另外还有很多可供选择的指数基金，它们又该如何去估值呢？我们来学习另一位大师的估值策略，约翰·博格使用的博格公式法。

\subsection{盈利快速增长：博格公式法}
如果一个品种长期盈利增长速度远高于10\%，那就不适合用盈利收益率法了。

对于盈利快速增长的品种，我们更多会考虑博格公式法。

约翰·博格是世界上第一只指数基金的发明者，被称为“指数基金之父”。约翰·博格也是一位成功的企业家，是世界第二大基金管理公司——先锋领航集团的创始人，掌管上万亿美元资金的运作。约翰·博格还与巴菲特、格雷厄姆和彼得·林奇（Peter Lynch）被财富杂志评为“20世纪四大投资巨人”。

博格公式分析了影响指数基金收益的几个因素，依靠博格公式，我们可以投资盈利高速增长的品种，或者盈利呈周期性变化的品种。

{\heiti 影响指数基金收益的几个因素}

{\kaishu 指数基金净值=市盈率×盈利+分红}

从上述公式可以看出，有三个因素直接影响了我们投资指数基金的收益：第一个是初始投资时刻指数基金的股息率，它影响了我们的分红收益；第二个是投资期内指数基金的市盈率变化；第三个是投资期内指数基金的盈利增长率。最后这两项影响了我们的资本利得收益。分红收益+资本利得收益，也就是我们投资的全部收益了。

股息率的获取非常容易。因为初始投资时的股息率，在我们买入指数基金时就确定了。一般来说，指数基金越是低估，也就是当其价格越低于其内在价值时，股息率越高。

主要是市盈率的变化和未来盈利的增长，判断起来相对困难。

{\heiti 市盈率}

{\kaishu 市盈率=公司市值/公司盈利}

对一个宽基指数来说，市盈率通常是在一个范围内来回波动，不会太高也不会太低。

比如说沪深300指数，历史上市盈率最高是48倍，最低是8倍左右。

熊市最底部的区域，沪深300指数的市盈率也处于低位，大约在8～11倍。牛市最顶部的区域，沪深300指数的市盈率也会大涨，大约在17～48倍。其他时间里，沪深300指数的市盈率，基本在8～17倍来回波动。

那么很明显，如果我们买入时的市盈率处于低位，之后市盈率上升，我们就可以从市盈率的变化上获得正收益。但是如果买入时的市盈率处于高位，之后市盈率下降，我们就会面临巨大的亏损。

例如2014年的时候，沪深300指数的市盈率在8～9倍。之后2014—2015年出现了一波牛市，沪深300指数上涨了两倍多，收益非常丰厚。

但是如果在2007年牛市，沪深300指数市盈率高达48倍的时候买入，之后沪深300指数基金短期里会出现“腰斩”，投资者亏损50\%以上。

这就是买入估值跟收益的关系。

{\heiti 盈利的长期增长和分红的积累}

但是能不能说，我们的收益是由估值变化带来的呢？

短期内，指数基金的涨跌，主要是由估值的变化决定的。但是时间拉长后，指数基金的上涨，是由背后公司的盈利增长和分红积累推动的。

因为估值的变化是有极限的，它只会在一个范围内来回波动，不具备复利增长性。指数基金真正具备复利增长能力的，是盈利的长期增长。

长期来看，对沪深300指数来说，盈利以平均9\%～10\%/年的速度增长。

分红，我们在前面介绍过，是盈利的一部分，所以盈利的增长和分红的积累，都是依靠背后公司的经营。

经济发展速度越快，上市公司的盈利增长速度也会越快。这是我们长期投资指数基金获得盈利的核心所在。

所以巴菲特说过，买指数就是买国运。得益于国内的经济高速发展，过去几十年，国内很多指数背后公司的盈利增速都是不错的。

而且时间越长，盈利增长带来的收益，占比会越高。

通常来说，在一轮牛熊市里（7年左右时间），估值变化带来的收益占比比较高，但是在两三轮牛熊市里，一定是盈利增长带来的收益高。

所以如果给估值和盈利增长分一下优先级的话，那一定是在盈利增长良好，其次是在估值处于低位的时候买入。

所以我们也就得到了改良后的定投指数基金的方法。

（1）挑选长期盈利增长率良好的品种。

其实这就是我们前面介绍过的宽基指数基金和优秀行业指数基金。它们基本都是符合这个要求的。

（2）在估值处于低位的阶段买入。

这个低位，要同时满足两个要求：一是在自身历史估值处于低位，二是估值要跟自身盈利增长速度匹配。

在自身历史估值处于低位时买入，未来回归正常估值区域时，我们就获得了估值从低变高的收益。反过来，如果估值处于高位，未来估值从高变低，会降低我们的投资回报。过低或者过高的估值，之后都会回归。

但仅仅看历史是不够的，估值也要跟自身盈利增长速度匹配。比如说美国股票市场的指数，长期盈利增长率大约是10\%，美国股市过去百年的市盈率平均值，大约是15倍。这就是一个对应关系。

假如说国内的某个指数，背后公司长期盈利增长率是15\%，这个指数能享受到的长期合理市盈率，也会更高一些。但是如果长期盈利增长速度只有6\%，它的市盈率还有二三十倍，那就不合理，两者并不匹配。

实际上，因为国内股票市场不成熟，过去很长时间内，存在短期炒作一些业绩很差的股票，或者只有题材和概念的股票。这些股票长期盈利增长率是个位数，甚至是负数，但市盈率经常被炒到几十倍甚至上百倍。

对于这类品种，可能20倍市盈率就是它们历史最低的估值区域，但仍然跟盈利增长率不匹配。所以，不能仅仅看估值处于历史低位，就盲目地认为这个品种是值得投资的。

这就是我们应用估值的方法。

从全世界范围来看，成熟国家的股票市场，长期盈利增长率在8\%～10\%，市盈率平均值在15倍左右。这是一个可以参考的估值区间。

如果盈利增长率更高，那可以享受的市盈率平均值也会更高。但是很少有品种能长期保持30\%以上的盈利增长速度，所以对市盈率数值超过30倍的指数基金，一定要慎之又慎。

{\heiti 估值和盈利增长速度的关系：PEG指标}

其实估值和盈利增长速度之间的关系，另一位投资大师彼得·林奇也提到过。彼得·林奇最出名的估值指标就是PEG指标。

{\kaishu PEG=市盈率数值/长期盈利增长速度数值}

如果一个品种，市盈率是15倍，长期盈利增长率是15\%，那PEG=1。

通常来说，PEG的合理区间是1～2，平均值也是1.5。这也跟成熟市场的估值和盈利增长速度相匹配。如果PEG小于1，也就是市盈率的数值小于未来长期盈利增长速度的数值，那还是比较好的。

更多时候，PEG是帮助我们排雷用的，减少风险。

PEG低的品种，并不一定是有投资价值的，需要具体看这个品种的盈利增长速度能不能长期持续。而PEG高的品种，必须要谨慎。因为很少有长期保持高盈利增长速度的品种，所以对任何市盈率数值过高的品种都需要格外小心，宁可不投资。

{\heiti 博格公式法}

综上所述，博格公式法可总结如下。

\begin{itemize}
\item 长期盈利增长速度大幅高于10\%的品种，适合使用博格公式法。
\item PEG反映了市盈率波动区间和长期平均盈利增长速度之间的关系，PEG通常在1～2。如果市盈率长期保持很高的数值，那么对这类品种要谨慎，因为很少有品种能长期保持较高盈利增长速度。
\item 在这个基础上，结合指数的历史市盈率，看看当前市盈率所处历史市盈率的位置。如果市盈率同时也处于低位，那这个品种大概率就是处于低估区域的，可以考虑买入。
\item 如果市盈率处于高位，或者PEG也大于2，那这个品种大概率就是处于高估区域，这时就要谨慎一些。
\end{itemize}

\subsection{改进后的定投法：收益更好、风险更低}
{\heiti 盈利收益率和博格公式的关系}

其实，盈利收益率法和博格公式的原理是相通的。

盈利收益率是市盈率的倒数。当盈利收益率大于10\%的时候，也就是市盈率小于10的时候。对绝大多数指数来说，市盈率小于10，都是处于历史市盈率波动范围的较低区域，意味着未来市盈率大概率是上涨的。

同时适合用盈利收益率法的品种，盈利增长比较稳定，也可以获得盈利增长的收益。这样估值从低变高获得了收益，盈利稳定增长获得了收益，两者叠加，也符合博格公式的原理。

换句话说，盈利收益率法，其实是博格公式的快速判断版本。

区别在于，盈利收益率法不太需要看历史估值数据，只需要对应品种的盈利长期比较稳定就可以使用。博格公式需要分析指数的历史估值区间和盈利增长情况。盈利增长良好的品种，自身估值处于历史较低区域，且估值数值跟盈利增长速度相匹配。

{\heiti 螺丝钉指数估值表}

看到这里，有的朋友可能还是会对估值有些迷惑。

估值确实是非常难的一件事情，哪怕是投资大师，估值也会出现错误。这短短的一章，也不可能就能把估值这件事情讲明白。只不过看完了这一章，我们能大致明白估值的几种方法。

指数的估值数据查找并不方便，并且使用盈利收益率法和博格公式法估值的过程也比较复杂。所以螺丝钉将各个指数的估值数据进行汇总，在每个交易日的晚上，公布当天收盘后各个指数基金的市盈率、市净率等估值。数据每个交易日更新，保证实效性，所以大家查询到的都是最新的数据信息。这个估值表已经连续更新一千多期，累计上亿人次观看，还是经受住了市场的考验的。并且这个估值表，应用了前面说的格雷厄姆的盈利收益率法（估值表里有盈利收益率的品种），以及博格公式法。

我们以2018年年初估值表为例，如图9.4所示，做一下讲解。

（1）左上角是估值表对应的日期。每个交易日晚上会进行更新。

（2）估值表用颜色代表低估高估。

用绿色（图9.4中第一部分）表示当前低估，这些指数基金是适合定投的。

用黄色（图9.4中第二部分）表示当前正常，如果之前有定投这些指数基金，可以暂停定投，继续持有。

用红色（图9.4中第三部分）表示当前高估，如果之前有持有这些指数基金，可以开始卖出。

（3）估值表后面有市盈率、市净率、股息率这三个最重要的估值指标。

（4）后面有两列指数代码，分别是场内和场外对应的指数基金。

（5）\$\textbackslash\{\}ast\$，代表不适合参考这个数据或者暂无。

我们在定投的时候，可以参考这个估值表，省去了自己估值的烦琐过程，直接应用到我们的定投中。每次选择绿色低估区域（如图9.4中第一部分）的指数基金定投，到了黄色正常估值区域（如图9.4中第二部分）持有，到了牛市，指数基金到红色高估区域（如图9.4中第三部分），就可以开始卖出。

\begin{figure}[htbp]
\centering
\includegraphics[width=\textwidth,height=0.82\textheight,keepaspectratio]{figure/Image00034.png}
\end{figure}
图9.4螺丝钉指数估值表

{\heiti 改进后的两大好处}

使用改进后的定投方法，我们能够享受到两大好处。

第一个好处，是能够获得更好的收益。买入价格低于价值的品种，我们可以赚取低估回归的额外收益，相比起普通的指数基金定投，收益要高出不少。

比如说我们回测一下，在恒生指数上，使用低估投资的方法。

历史上，恒生指数每次进入低估的时候开始买入，进入高估之后卖出。

第一次，1974—1975年，从170点涨到400点，年复合收益率达到200\%以上。

第二次，1982—1986年，从1200点涨到2300点，年复合收益率是17.6\%。

第三次，1989—1993年，从3000点涨到9000点，年复合收益率达到31\%。

第四次，1998—1999年，从9000点涨到14000点，年复合收益率达到50\%。

第五次，2008—2009年，从15000点涨到20955点，年复合收益率是39\%。

在这几次定投期间，投资者都能获取平均15\%以上的年复合收益率。

不过也有朋友会有疑问：恒生指数并不是每个时间段都有投资的机会，如果之间没有低估机会，岂不是只能白白等待？

其实不要紧的，对港股指数基金来说，通过对历史的分析，基本上每隔4～5年就会出现一次机会。这还只是港股，如果加上A股、美股、其他股票市场，平均每2～3年就有一次低估机会。

以最近10年为例，2008—2010年的美股、2012—2014年的A股、2016年年初的港股、2018年下半年的A股，都出现了很不错的投资机会。

再加上国内外众多的指数基金，不同指数基金的成长性是不同的，有的处在低估区域，有的处在正常区域，有的处在高估区域，所以我们不用担心买不到。所有指数基金都不低估，这种情况是很罕见的。

第二个好处，是投资风险更低。这是由于价格低于价值的品种，自身下跌空间就比较有限，投资的风险就更低一些。再加上定投这种持续、多次的投资方式，又可以进一步分摊风险，降低波动。定投低估值的品种，将指数基金的风险又大幅降低了。

估值并不是一件容易的事情，不过好在我们有很多投资大师的估值方法可以借鉴。

我们主要应用市盈率、市净率、股息率，以及背后公司的盈利增长速度这几个指标，结合格雷厄姆的盈利收益率法、约翰·博格的博格公式法、彼得·林奇的PEG指标，可以帮助我们解决大部分指数基金品种的估值问题。

这套估值方法，螺丝钉应用多年，也经受住了牛熊市的考验，有几十万投资者也在使用这套方法，还是很稳定可靠的。不过就算如此，我们应用估值方法，计算出来还是会有误差。我们要保留一定的安全边际，为自己留出容错的空间。

再配合定投的方式，把不必要的风险降到最低。定投的方法跟估值的方法相配合，可以得到“1+1>2”的效果。

掌握长期收益更好的优秀指数基金品种，然后耐心等待熊市低估的阶段定投它们，这是最适合普通投资者的投资技巧。

{\heiti 投资者笔记}

\begin{itemize}
\item {\kaishu 投资者投资指数基金最大的一个误区，就是认为一只指数基金好就自始至终去投资它，不管贵贱。}
\item {\kaishu 在价格低于价值时买入，才能买得划算。没有人愿意用1元买价值0.8元的东西。}
\item {\kaishu 常见的估值方法可以分为：相对估值法和绝对估值法。}
\item {\kaishu 常见的估值指标有：市盈率、市净率、股息率。}
\item {\kaishu 不同的指数基金需要使用不同的估值方法进行投资。比如盈利始终保持不变的品种可以使用现金流收益率进行判断；盈利稳定增长的品种可以使用盈利收益率法进行判断；盈利快速增长的品种可以使用博格公式法进行判断。}
\item {\kaishu 快速又方便查询指数估值数据的方式：通过螺丝钉的公众号获取每日估值数据推送。}
\end{itemize}

